\uuid{kk7a}
\titre{Maximum}
\auteur{bodin ; hanani ; mzari}
\organisation{exo7}
\licence{CC BY-NC-SA 4.0}
\datecreate{2020-01-20}
\niveau{L1}
\module{Analyse}
\chapitre{Propriétés de R}
\sousChapitre{Maximum, minimum, borne supérieure}
\difficulte{1}
\theme{Maximum, majorant, maximum}
\exo7qcmid{lille-1.reels.028}

\contenu{
\begin{question}
\qclassification{120.02}
\qtags{motcle=maximum}
Le maximum d'un ensemble $E$, s'il existe, est le réel $m \in E$ tel que pour tout $x\in E$, on a $x \le m$.
\begin{answers}
    \bad{Si $E = [3,7]$ alors $8$ est un maximum de $E$.}

    \good{Si $E = [-3,-1]$ alors $-1$ est le maximum de $E$.}

    \good{L'ensemble $E = [-3,-1[$ n'admet pas de maximum.}

    \bad{L'ensemble $E = [-3,2[ \  \cap \  ]-1,1]$ n'admet pas de maximum.}

\end{answers}
\begin{explanations}
Attention, le maximum de $E$ doit être un élément de $E$ !
\end{explanations}
\end{question}
}
